\documentclass[10pt,a4paper]{article}
\usepackage[utf8]{inputenc}
\usepackage[T1]{fontenc}
\usepackage[french]{babel}
\usepackage{geometry}
\geometry{left=1cm, right=1cm, top=1cm, bottom=1cm}
\usepackage{enumitem}
\usepackage{hyperref}
\usepackage{fontawesome5}
\usepackage{multicol}
\usepackage{graphicx}
\setlength{\columnsep}{0.6cm}
\setlist{nosep, left=0pt, itemsep=0pt, topsep=0pt, partopsep=0pt, parsep=0pt}
\pagestyle{empty}

\begin{document}

% ========== HEADER ==========
\begin{minipage}{0.2\textwidth}
    \includegraphics[width=\linewidth]{photo.jpg}
\end{minipage}
\begin{minipage}{0.6\textwidth}
    \centering
    {\huge \textbf{Julien GONZALES}} \\[0.2cm]
    \textbf{Apprenti Ingénieur en Systèmes Électriques et Électroniques Embarqués}\\
    \textit{Recherche contrat d'apprentissage de 36 mois}\\[0.1cm]
\end{minipage}
\begin{minipage}{0.2\textwidth}
    \raggedleft
    
\includegraphics[width=2.5cm]{logo_cefipa.png}\\[0.2cm]
    
\includegraphics[width=2.5cm]{logo_cesi.png}
\end{minipage}

\vspace{0.2cm}

% ========== CONTACT INFO ==========
\noindent
\faEnvelope\ \href{mailto:julienabdougonzales@gmail.com}{julienabdougonzales@gmail.com} \quad
\faPhone\ 07.75.85.70.82 \quad
\faMapMarker\ 91400 Orsay / Mobilité France entière \quad
\faCar\ Permis B

\vspace{0.3cm}

\begin{multicols}{2}

\section*{Objectif}
Intégrer Naval Group en alternance comme Ingénieur Automaticien pour concevoir et valider des systèmes de contrôle embarqués.

% ========== FORMATION ==========
\section*{Formation}
\textbf{Programme ingénieur S3E (Bac+5) – CESI Nanterre} \hfill \textit{2026 – 2029} \\
Systèmes embarqués, IA temps réel, FPGA, RTOS.

\textbf{Cursus génie électrique (Bac+3) – Johns Hopkins University} \hfill \textit{2022 – 2025} \\
Mention très bien – GPA 3,44/4,0. Spécialisation : systèmes embarqués, traitement du signal, capteurs.


% ========== EXPÉRIENCES PROFESSIONNELLES ==========
\section*{Expériences professionnelles}
\textbf{Câbleur – Electro Faisceaux} \hfill \textit{depuis janvier 2026}
\begin{itemize}[leftmargin=*]
    \item \textbf{Assemble et intègre} des faisceaux électriques complexes pour systèmes industriels et dispositifs médicaux critiques.
    \item \textbf{Diagnostique, corrige et valide} les défauts électriques en environnement normé (traçabilité, qualité, conformité).
\end{itemize}

\textbf{Research Intern – RISE@APL (Johns Hopkins APL)} \hfill \textit{2024 – 2025}
\begin{itemize}[leftmargin=*]
    \item Concevoir des systèmes automatisés selon cahier des charges techniques.
    \item Développer des logiciels automates et créer schémas électriques embarqués.
    \item Réaliser et valider les essais techniques jusqu'à mise en service.
\end{itemize}

\textbf{Assistant laboratoire – Johns Hopkins University} \hfill \textit{2024 – 2025}
\begin{itemize}[leftmargin=*]
    \item \textbf{Conçoit et exécute} des protocoles de validation pour systèmes électroniques de précision.
    \item \textbf{Structure et formalise} la documentation technique pour reproductibilité expérimentale.
\end{itemize}

% ========== PROJETS ==========
\section*{Projets – Réalisations concrètes}
\textbf{CNN pour reconnaissance d’images – TensorFlow/PyTorch} \hfill \textit{2024 – 2025}
\begin{itemize}[leftmargin=*]
    \item \textbf{Développe et entraîne} des architectures CNN pour classification (CIFAR-10, MNIST).
    \item \textbf{Réduit la latence et l’empreinte mémoire} via quantification et optimisation embarquée.
\end{itemize}

\textbf{Ondelettes sur FPGA (DWT)} \hfill \textit{2024}
\begin{itemize}[leftmargin=*]
    \item \textbf{Implémente et pipeline} une DWT sur FPGA (VHDL) pour traitement temps réel.
    \item \textbf{Démontre un gain de performance} sur le débruitage d’images infrarouges.
\end{itemize}

\textbf{Drone – Fusion de capteurs \& PID} \hfill \textit{2024 – 2025}
\begin{itemize}[leftmargin=*]
    \item \textbf{Développe un firmware embarqué robuste} (STM32, C) intégrant IMU (MPU6050) via I2C.
    \item \textbf{Implémente une fusion de capteurs temps réel} et une boucle PID stable (<10 ms).
\end{itemize}

\textbf{Oxymètre – Pipeline Linux embarqué} \hfill \textit{2024}
\begin{itemize}[leftmargin=*]
    \item \textbf{Construit une chaîne complète} d’acquisition et traitement (SPI, filtrage, calcul SpO2).
    \item \textbf{Industrialise les tests} via scripts Python/Bash et gestion de version Git.
\end{itemize}

% ========== RECHERCHE – CSMS LAB ==========
\section*{Recherche – CSMS Lab (Johns Hopkins)}
\textit{3 ans – apprentissage alternatif, systèmes haute dimension}
\begin{itemize}[leftmargin=*]
    \item \textbf{Conçoit des algorithmes d’apprentissage locaux} comme alternative scalable à la rétropropagation.
    \item \textbf{Explore les propriétés des systèmes en grande dimension}.
    \item \textbf{Relie théorie et systèmes physiques} (neuromorphique, calcul distribué, capteurs intelligents).
\end{itemize}

\end{multicols}

\vspace{0.1cm}
\section*{Compétences}
\textbf{Langages} : C, C++, Python, VHDL, Bash, Matlab \\
\textbf{Systèmes embarqués} : STM32, ESP32, ATmega32, FPGA, Automates Schneider/Siemens, RTOS \\
\textbf{Outils} : Git, SEE Electrical Expert, Cadence, Kicad, Oscilloscope
\vspace{0.2cm}

% ========== LANGUES & CENTRES D'INTÉRÊT ==========
\noindent
\textbf{Langues} : Français (natif), Anglais (bilingue – TOEIC 950) \hfill
\textbf{Centres d’intérêt} : Course à pied, photographie, veille technologique.

\end{document}